\scanpage{63}\mspage{10}
une application canonique
\[
\Sigma_{\underline J}\xrightarrow{\ f^{t}_{\unsure{2|2_0}}\ }\Sigma_{\underline J}
\tag*{(30)}
\]
\margin{c'est l'application déduite par passage aux quotients de $2\operatorname{id}_{E_0}:E_0\longrightarrow E_0$}
% slanted
\margin{il faudrait revenir sur une description intrinsèque de $f_{\unsure{i|2}}^t$, sans faire intervenir \unsure{de courbe elliptique|des courbes elliptiques|de courbes elliptiques}}
qui doit pouvoir se définir en termes de $\underline J$, et de la seule section $t$ de $\Sigma_{\underline J}$ (ce qui en effet détermine $E_0$, à isom \inserted{(non unique)} près loc. pour la top. étale\ldots) On a ici
\[
\deleted{I=J\amalg}\qquad \underline I_t=\underline J\amalg t(S)
\tag*{(31)}
\]
\margin{\uline{NB} comme $S$-schéma, $\underline I_t$ ne dépend pas de $t$, mod iso. can.}
et \deleted{$\unsure{i|I|E_0}$} $\underline V$ peut se décrire comme $V_{\underline I}\simeq V_{\underline J}$, opérant sur $\underline I$ de la façon canonique, cette opération se prolongeant en une opération sur $\Sigma_{\underline J}$ (dépendant de la section $t$ de $\Sigma_{\underline J}^*$). Le \deleted{applic} morphisme $f_t$ est défini sans ambiguïté par les conditions suivantes (où on suppose pour simplifier que $\underline J=J_S$, $J$ ens à trois éléments) :
\[
\left\{\begin{aligned}
&\forall i\in J,\ \text{si }\sigma_i^t\in\operatorname{Aut}_S(\Sigma_J)
  \text{ est l'autom. qui échange}\\
&\qquad t\text{ et }i,\text{ et }j\text{ et }k\quad(\text{où }\{i,j,k\}=J),\text{ et}\\
&\deleted{\text{alors }\unsure{\text{si}|\text{soit}|f}\,(\Sigma_J)^{\sigma_i^t}}
  \qquad\Delta_i^t=(\Sigma_J)^{\sigma_i^t},\ \text{on a}\\
&\hspace{8em}f(\Delta_i^t)\subset\text{section }\zeta_i\text{ de }D_{\underline J}\simeq J_S
\end{aligned}\right.
\tag*{(32)}
\]
et cela doit impliquer automatiquement
\[
f_t(I_S=J_S\cup t(S))\subset\deleted{\text{section}}\;t(S)
\tag*{(33)}
\]
% archive sheet number: 62
